\documentclass[10pt,a4paper]{amsart}
\usepackage[margin=19mm]{geometry}
\usepackage{amsmath,amssymb,mathtools}
\usepackage[scr=rsfs,cal=euler]{mathalfa}
\usepackage{xfrac}
\usepackage[unicode,hidelinks,bookmarks=false]{hyperref}
\newcommand{\ie}{i.e.\ }
\newcommand{\cf}{cf.\ }
\newcommand{\resp}{respectively\ }
\newcommand{\Subsection}{\subsection}
\providecommand{\nobreakdash}{\nobreak\hbox{-}}
\setlength{\emergencystretch}{2em}
\multlinegap0pt
\usepackage{amssymb}
\usepackage{xcolor}
\usepackage{algorithm}\usepackage{algpseudocode}
\newtheorem{proposition}{Proposition}
\title{Context Channel Capacity: An Information-Theoretic Framework\\for Understanding Catastrophic Forgetting}
\author{Ran Cheng}
\date{Source: arXiv:2603.07415v1 (CC BY 4.0)}
\begin{document}
\maketitle

\noindent\emph{Source and licence.} Context Channel Capacity: An Information-Theoretic Framework\\for Understanding Catastrophic Forgetting. Version arXiv:2603.07415v1 (CC BY 4.0), distributed by arXiv under the Creative Commons Attribution 4.0 International licence, http://creativecommons.org/licenses/by/4.0/, as recorded in the arXiv licence metadata for this exact version and shown on the abstract page https://arxiv.org/abs/2603.07415v1. Full licence notice: \texttt{licenses/arxiv-2603.07415-CC-BY-4.0.txt}.\par\smallskip

\noindent\textit{Editorial scope.} Counted unit: the article's proposition prop:bypass (source0.tex lines 1449--1452). The article's complete original proof of it is delivered with it (source0.tex lines 1454--1497, one \texttt{proof} environment of 44 lines). No mathematics is written, repaired or re-derived: the statement and the proof are the article's own lines, in the article's own notation, and no retained line is substituted or re-flowed.  No further source-local context is needed: every cross-reference of every retained line resolves inside the two delivered spans.  Nothing was omitted from any retained span, so the package carries no omission record.  No retained line contains a citation, so the wrapper carries no bibliography and no bibliography entry.  No retained line contains a figure environment, an image-inclusion command or drawing code, so no image is delivered and the package declares no asset dependency.  Editorial formatting only, in addition: an AMS \texttt{amsart} preamble that restates the article's own declarations and macros used by the retained lines, namely the environments \texttt{proposition} on separate counters, together with the standard packages those lines need.  The retained environments print in this document as Proposition 1 in order of appearance, whereas the article prints the same items with the numbering of its own full document.  The complete native source of the article as distributed in the arXiv e-print for this exact version is bundled unchanged as \texttt{sources/195942/source0.tex}.
\begin{proposition}[Base Parameter Bypass in HyperNetworks]
\label{prop:bypass}
In a meta-learning framework where $\theta_k = \theta_\mathrm{base} + \mathbf{V} g_\psi(c_k)$ and $\theta_\mathrm{base}$ is a free parameter, the gradient signal through the base path dominates the context path when $\mathbf{V}$ is initialized with small scale, causing ``bypass collapse'' where $\theta_\mathrm{base}$ absorbs all task information.
\end{proposition}

\begin{proof}
\textbf{Step 1: Gradient magnitudes.}
Consider the meta-loss:
\begin{equation}
    \mathcal{L}(\theta_\mathrm{base}, \mathbf{V}, \psi) = \sum_{k=1}^K \mathbb{E}_{(x,y) \sim D_k}\left[\ell\bigl(f(x;\, \theta_\mathrm{base} + \mathbf{V} g_\psi(c_k)),\; y\bigr)\right]
\end{equation}
Let $\theta_k = \theta_\mathrm{base} + \mathbf{V} g_\psi(c_k)$. The gradients with respect to each parameter group are:

\emph{Base path:}
\begin{equation}
    \frac{\partial \mathcal{L}}{\partial \theta_\mathrm{base}} = \sum_{k=1}^K \mathbb{E}\left[\frac{\partial \ell}{\partial \theta_k}\right]
\end{equation}
This is simply the sum of per-task gradients, with no additional Jacobian factors.

\emph{Context path (through $\psi$):}
\begin{equation}
    \frac{\partial \mathcal{L}}{\partial \psi} = \sum_{k=1}^K \mathbb{E}\left[\frac{\partial \ell}{\partial \theta_k} \cdot \mathbf{V} \cdot \frac{\partial g_\psi(c_k)}{\partial \psi}\right]
\end{equation}
This involves the additional factor $\mathbf{V} \cdot \frac{\partial g_\psi}{\partial \psi}$.

\textbf{Step 2: Scale analysis.}
At initialization with $\mathbf{V} \sim \mathcal{N}(0, \sigma_V^2)$ where $\sigma_V = 0.01$:
\begin{equation}
    \left\|\frac{\partial \mathcal{L}}{\partial \theta_\mathrm{base}}\right\| = O\left(\left\|\frac{\partial \ell}{\partial \theta}\right\|\right), \qquad
    \left\|\frac{\partial \mathcal{L}}{\partial \psi}\right\| = O\left(\sigma_V \cdot \left\|\frac{\partial \ell}{\partial \theta}\right\| \cdot \left\|\frac{\partial g}{\partial \psi}\right\|\right)
\end{equation}
The ratio is:
\begin{equation}
    \frac{\|\partial \mathcal{L} / \partial \psi\|}{\|\partial \mathcal{L} / \partial \theta_\mathrm{base}\|} = O\left(\sigma_V \cdot \left\|\frac{\partial g}{\partial \psi}\right\|\right) \approx O(0.01)
\end{equation}
assuming $\|\partial g / \partial \psi\| = O(1)$ at initialization.

\textbf{Step 3: Information absorption.}
Since the base path gradient is $\sim 100\times$ larger, Adam (or any adaptive optimizer) moves $\theta_\mathrm{base}$ much faster than $\psi$. After a few optimization steps, $\theta_\mathrm{base}$ has absorbed the ``average task'' solution, and the per-task residuals $\ell(f(x; \theta_\mathrm{base}), y)$ are small. This further reduces $\|\partial \ell / \partial \theta\|$, making the context path gradient even smaller---a positive feedback loop that leads to bypass collapse: $g_\psi(c_k) \approx \mathrm{const}$ for all $k$.

\textbf{Step 4: Base dropout as a solution.}
Setting $\theta_\mathrm{base} \to 0$ with probability $p$ (base\_drop\_rate) during training removes the base path for a fraction $p$ of training steps, forcing:
\begin{equation}
    \theta_k = \mathbf{V} g_\psi(c_k) \cdot \frac{\|\theta_\mathrm{base}\|}{\|\mathbf{V} g_\psi(c_k)\| + \epsilon}
\end{equation}
(with appropriate rescaling to maintain output magnitude). On these steps, $\partial \mathcal{L}/\partial \psi$ is the \emph{only} gradient pathway, forcing the context encoder to learn task-discriminative representations. Our experiments confirm that $p = 0.3$ is sufficient: P5 delta drops from $0.0$ (collapsed) to $-96.4$ percentage points (fully context-dependent) on Split-MNIST.

\emph{Critical note:} Base dropout alone is insufficient if the context encoder uses non-differentiable statistics (e.g., EMA running mean/variance). The \texttt{direct\_stats=True} fix (using current-batch statistics directly, which are differentiable) is a necessary prerequisite.
\end{proof}

\end{document}
